\section{Dataset}
\label{sec:dataset}

\subsection{Source Material}

ECAG-Bench is derived from \emph{Examples and Counterexamples in Elementary Algebraic Geometry}, a collection of worked examples originally written as supplementary material for learners entering algebraic geometry.
The collection has been developed over several years and covers topics ranging from basic scheme theory to derived categories and moduli spaces.

The original material was written in \LaTeX\ and subsequently converted to Markdown for web publication.
Each example was then enhanced to meet a uniform quality standard: a self-contained mathematical narrative with complete constructions, precise references, and explicit computational details.

\subsection{Content Enhancement}

The content enhancement process involved restructuring each example into a natural mathematical exposition following a consistent quality standard:
\begin{enumerate}[nosep]
  \item \textbf{Narrative form:} Each example reads as a self-contained mathematical essay rather than a template with bullet points.
  \item \textbf{Complete constructions:} All computational steps are shown explicitly, allowing independent verification.
  \item \textbf{Benchmark problem extraction:} Each example is annotated with a structured benchmark problem specification, including the problem statement, expected answer type, difficulty level, and verification strategy.
\end{enumerate}

All 332 entries have been enhanced to this standard, with zero stubs remaining.

\subsection{Statistics}

\begin{table}[h]
\centering
\caption{ECAG-Bench dataset statistics.}
\label{tab:statistics}
\begin{tabular}{lr}
\toprule
\textbf{Metric} & \textbf{Count} \\
\midrule
Total entries & 332 \\
Substantive (with solution) & 332 \\
Stubs & 0 \\
\midrule
\multicolumn{2}{l}{\textbf{By Chapter}} \\
\midrule
Basic Concepts & 162 \\
Computations & 66 \\
Moduli \& Deformation & 58 \\
Cohomology & 46 \\
\midrule
\multicolumn{2}{l}{\textbf{By Difficulty}} \\
\midrule
Level 1 (Routine) & 77 \\
Level 2 (Non-trivial theorem) & 237 \\
Level 3 (Deep/combining results) & 18 \\
\midrule
\multicolumn{2}{l}{\textbf{By Question Type}} \\
\midrule
Example & 243 \\
Counterexample & 66 \\
Computation & 14 \\
Proof & 9 \\
\midrule
\multicolumn{2}{l}{\textbf{By Verification Strategy}} \\
\midrule
CAS & 309 \\
Exact match & 14 \\
Rubric & 9 \\
\bottomrule
\end{tabular}
\end{table}

\subsection{Difficulty Levels}

Problems are assigned one of three difficulty levels:
\begin{description}[nosep]
  \item[Level 1 (Routine):] The construction follows directly from definitions. A student who has read the relevant section of a standard textbook should be able to produce an answer.
  \item[Level 2 (Non-trivial):] The construction requires applying a significant theorem or combining multiple standard results. This is the level of a good homework problem in a graduate course.
  \item[Level 3 (Deep):] The construction requires either deep structural insight, the combination of results from different areas, or the knowledge of a specific ``trick'' that is not obvious from the problem statement.
\end{description}

\subsection{Quality Control}

Each entry has been reviewed for:
\begin{itemize}[nosep]
  \item Mathematical correctness of the stated construction
  \item Completeness of computational details
  \item Accuracy of the benchmark problem formulation
  \item Consistency of difficulty rating across the dataset
  \item Validity of the assigned verification strategy
\end{itemize}
